\documentclass[10pt,a4paper]{amsart}
\usepackage[margin=19mm]{geometry}
\usepackage{amsmath,amssymb,mathtools}
\usepackage[scr=rsfs,cal=euler]{mathalfa}
\usepackage{xfrac}
\usepackage[unicode,hidelinks,bookmarks=false]{hyperref}
\newcommand{\ie}{i.e.\ }
\newcommand{\cf}{cf.\ }
\newcommand{\resp}{respectively\ }
\newcommand{\Subsection}{\subsection}
\providecommand{\nobreakdash}{\nobreak\hbox{-}}
\setlength{\emergencystretch}{2em}
\multlinegap0pt
\usepackage{bm}
\usepackage{bbm}
\usepackage{dsfont}
\usepackage{xspace}
\usepackage{cancel}
\usepackage{mathtools}
\usepackage{amsthm}
\makeatletter
\@ifundefined{theorem}{\newtheorem{theorem}{Theorem}}{}
\@ifundefined{lemma}{\newtheorem{lemma}[theorem]{Lemma}}{}
\makeatother
\title[Weak solutions and incompressible limit of a quasi-incompres]{Weak solutions and incompressible limit of a quasi-incompressible Navier--Stokes/Cahn--Hilliard model for viscous two-phase flows}
\author{Mingwen Fei, Xiang Fei, Daozhi Han, Yadong Liu}
\date{Source: arXiv:2508.08090v1 (CC BY 4.0)}
\begin{document}
\maketitle

\noindent\emph{Source and licence.} Weak solutions and incompressible limit of a quasi-incompressible Navier--Stokes/Cahn--Hilliard model for viscous two-phase flows. Version arXiv:2508.08090v1 (CC BY 4.0), distributed under the Creative Commons Attribution 4.0 International licence, as recorded in the arXiv licence metadata for this exact version and shown on the abstract page https://arxiv.org/abs/2508.08090v1. Full rights notice: licenses/arxiv-2508.08090-CC-BY-4.0.txt.\par\smallskip

\noindent\textit{Editorial scope.} Counted unit: the Lemma whose source label is \texttt{lem:higher eatimate} in the article (source0.tex lines 1315--1323, source label \texttt{lem:higher eatimate}), delivered together with the article's own complete original proof of it (source0.tex lines 1324--1606, one \texttt{proof} environment of 283 lines). No mathematics is written, repaired or re-derived: statement and proof are the article's own text line for line. Required context retained uncounted: model4-3 (lines 610--627); weak1 (lines 664--669); ConEnL2 (lines 683--687). Every definition, display and label of those lines is retained unchanged. Omitted from this package and not redistributed: 1--609, 628--663, 670--682, 688--1314, 1607--2926. Nothing inside a retained excerpt is omitted or altered: every retained body is delivered exactly as the article's lines are written, so no omission receipt of any kind accompanies this package. Citations: the retained excerpts carry no citation command at all, so no bibliography entry of the article is transcribed and no reference is redistributed. Reference adaptation: none was needed and none was made; every reference inside the delivered unit resolves inside it, so no unresolved reference or citation is produced. Printed slips kept literal: no printed word, symbol, hypothesis or number of the counted statement or of its proof was altered. Layout and class: the article's own class is \texttt{amsart} with packages including amsmath, amsthm, amssymb, mathrsfs, geometry, graphicx, mathtools, xcolor, todonotes, tikz, bm, stmaryrd, enumitem, appendix; the wrapper is the lane's pinned \texttt{amsart} editorial wrapper at 10pt on A4, which restates the theorem-like environments the retained text needs and adds the article's own macro definitions that the retained text needs (none), with extra wrapper packages (bm, bbm, dsfont, xspace, cancel, mathtools). The delivered numbering is the unit's own. Assets: no retained span contains a figure environment, a graphics-inclusion command, a PDF-page inclusion, a table or a TikZ picture; no image file is copied into this package, the package declares no asset dependency and no image file is redistributed with it. Licence: version arXiv:2508.08090v1 of this article is distributed under the Creative Commons Attribution 4.0 International licence, as recorded in the arXiv licence metadata for this exact version and shown on the abstract page https://arxiv.org/abs/2508.08090v1. A full rights notice is delivered at licenses/arxiv-2508.08090-CC-BY-4.0.txt. Characters: every retained excerpt is pure ASCII, so the delivered engine, XeLaTeX with its default Latin Modern text font, reports no missing character.
	\begin{align}
		\label{model4-1}
		\partial_t(\rho  \mathbf{u}) + \mathrm{div}\left (\rho  \mathbf{u}\otimes
		\mathbf{u}\right) 
		& = \mathrm{div}\big(S(\phi,\mathbb{ D}\mathbf{u})\big)
		- \zeta\rho\nabla p_{0}
		- \phi \nabla \mu_{p}^{0}
		+ \frac{\varepsilon\delta}{4\alpha}  p_{0} \mathbf{u}, 
		\\\label{model4-2}
		\partial_t \rho+\mathrm{div} (\rho \mathbf{u})& =\frac{\varepsilon\delta}{2\alpha} p_{0}, 
		\\
		\label{model4-3}
		\partial_t\phi + \mathrm{div}(\phi  \mathbf{u})& =\Delta\mu_p^{0}, \\
		\label{model4-4}
		\mu& =F'(\phi)+\Lambda^{2s} \phi, \\
		\label{model4-5}
		\mu_p^{0}+\overline{\mu}& =\mu+\alpha p_{0}, \\\rho& =\frac{\varepsilon}{2}\phi+\frac{\varepsilon}{2}+1,
	\end{align}

		\begin{align}\label{weak1}
			&\nonumber(\mathbf{u},\partial_{t}\boldsymbol{\varphi})_{Q_{T}}+(\mathbf{u}\cdot\nabla\mathbf{u},\boldsymbol{\varphi})_{Q_{T}}-(\rho^{-1}S(\phi,\mathbb{ D}\mathbf{u}),\nabla\boldsymbol{\varphi})_{Q_{T}}-(S(\phi,\mathbb{ D}\mathbf{u}),\nabla\rho^{-1}\otimes\boldsymbol{\varphi})_{Q_{T}}
			\\&\qquad
			+\zeta(p_{0},\mathrm{div}\boldsymbol{\varphi})_{Q_{T}}-(\rho^{-1}\phi \nabla \mu_{p}^{0},\boldsymbol{\varphi})_{Q_{T}}
			-(\frac{\varepsilon\delta}{4\alpha}  p_{0} \mathbf{u},\boldsymbol{\varphi}\rho^{-1})_{Q_{T}}=0,
		\end{align}

		\begin{align}\label{ConEnL2}
			\nonumber E(\phi(t),\mathbf{u}(t))&+\int_0^{t}\int_{\mathbb{T}^{3}}S(\phi,\mathbb{ D}\mathbf{u}):\mathbb{ D}\mathbf{u}\,\mathrm d x\mathrm dt
			\\&
			+\int_0^{t}\int_{\mathbb{T}^{3}}|\nabla \mu_{p}^{0}|^{2}\,\mathrm d x\mathrm dt+\delta\int_0^{t}\int_{\mathbb{T}^{3}}|{p}_{0}|^{2}\,\mathrm d x\mathrm dt\leq E(\phi_0,\mathbf{u}_0),
		\end{align}

			\begin{lemma}\label{lem:higher eatimate}
				
				Let $s>\frac{3}{2}$, then it follows
				\begin{align*}
					\|\phi_\delta\|_{L^{2}(0,T;H^{s+\frac{\gamma}{2}}(\mathbb{T}^{3}))}
					\leq C
				\end{align*}
				for  $0\leq\gamma\leq 1$, where $C>0$ depends on the initial data and $ T> 0$, but is independent of $\delta$ and $\alpha$.
			\end{lemma}

			\begin{proof}
				
				It follows from \eqref{ConEnL2} that
				\begin{align*}
					&\mathbf{u}_\delta\cdot\nabla\mathbf{u}_\delta, p_{0,\delta} \mathbf{u}_\delta\in L^\frac{4}{3}(0,T;L^{\frac{6}{5}}(\mathbb{T}^3)),
					\\& \rho_\delta^{-1}S(\phi_{\delta},\mathbb{ D}\mathbf{u}_\delta),\rho_\delta^{-1}\phi_\delta \nabla \mu_{p,\delta}^{0}\in L^2(0,T;L^{2}(\mathbb{T}^3)),
					\\& \nabla\rho_\delta^{-1}\cdot S(\phi_\delta,\mathbb{ D}\mathbf{u}_\delta)\in L^2(0,T;L^{\frac{6}{5}}(\mathbb{T}^3)).
					%\\&\rho_\delta^{-1}\phi_\delta \nabla \mu_{p,\delta}^{0} \in L^2(0,T;L^{2}(\mathbb{T}^3)),
					%\\&\frac{\varepsilon\delta}{4\alpha}  p_{0,\delta} \mathbf{u}_\delta\in L^\frac{4}{3}(0,T;L^{\frac{6}{5}}(\mathbb{T}^3)).
				\end{align*}
				Therefore  \eqref{weak1} holds for  $\boldsymbol{\varphi}\in  H^1(0,T;H^{-1}(\mathbb{T}^3))\cap L^4(0,T;L^{6}(\mathbb{T}^3))\cap L^2(0,T;H^{1}(\mathbb{T}^3))$.
				
				% For $(\mathbf{u}^{N},p_{0}^{N},\phi^{N},\mu_{p}^{0,N})$ in Section \ref{sec:approximation}, by utilizing \eqref{im1-w} and  integration by parts one has
				%\begin{align}\label{weak1110}
				%  &\nonumber(\mathrm{div}\mathbf{u}^{N},\partial_{t,h}^{+}(\psi\Delta^{-1}\Lambda^{\gamma}\phi_{\delta}))_{Q_{T}}+(\rho^{N}(\rho_{h}^{N})^{-1}\mathbf{u}^{N}\cdot\nabla \mathbf{u}^{N},\psi\nabla\Delta^{-1}\Lambda^{\gamma}\phi_{\delta})_{Q_{T}}
				%  \\&\nonumber\quad+(\nabla(\rho_{h}^{N})^{-1}S(\phi_{h}^{N},\mathbb{ D}\mathbf{u}^{N}),\psi\nabla\Delta^{-1}\Lambda^{\gamma}\phi_{\delta})_{Q_{T}}+((\rho_{h}^{N})^{-1}S(\phi_{h}^{N},\mathbb{ D}\mathbf{u}^{N}),\nabla(\psi\nabla\Delta^{-1}\Lambda^{\gamma}\phi_{\delta}))_{Q_{T}}
				% \\&\nonumber\quad -\zeta((\rho_{h}^{N})^{-1}\rho^{N}p_{0}^{N},\mathrm{div}(\psi\nabla\Delta^{-1}\Lambda^{\gamma}\phi_{\delta}))_{Q_{T}}
				%  +((\rho_{h}^{N})^{-1}\phi^{N}\nabla\mu_{p}^{0,N},\psi\nabla\Delta^{-1}\Lambda^{\gamma}\phi_{\delta})_{Q_{T}}
				%  \\&\quad-\zeta(\nabla((\rho_{h}^{N})^{-1}\rho^{N}) p_{0}^{N},\psi\nabla\Delta^{-1}\Lambda^{\gamma}\phi_{\delta})_{Q_{T}}
				% +(\frac{\varepsilon\delta}{4\alpha}(\rho_{h}^{N})^{-1}\mathbf{u}^{N}p_{0}^{N},\psi\nabla\Delta^{-1}\Lambda^{\gamma}\phi_{\delta})_{Q_{T}}=0,
				%\end{align}
				%where
				Due to the Sobolev embedding and elliptic regularity theory  we have
				\begin{align*}
					&\|\nabla\Delta^{-1}\Lambda^{\gamma}(\psi\phi_{\delta})\|_{L^4(0,T;L^{6}(\mathbb{T}^3))} \leq
					C\|\Lambda^{\gamma}(\psi\phi_{\delta})\|_{L^{\infty}(0,T;L^{2}(\mathbb{T}^{3}))}
					\leq C\|\psi\phi_{\delta}\|_{L^{\infty}(0,T;H^{\gamma}(\mathbb{T}^{3}))},
					\\&
					\|\nabla\Delta^{-1}\Lambda^{\gamma}(\psi\phi_{\delta})\|_{L^{2}(0,T;H^{1}(\mathbb{T}^{3}))} \leq
					C\|\Lambda^{\gamma}(\psi\phi_{\delta})\|_{L^{2}(0,T;L^{2}(\mathbb{T}^{3}))}
					\leq C\|\psi\phi_{\delta}\|_{L^{2}(0,T;H^{\gamma}(\mathbb{T}^{3}))},
					\\&
					\|\nabla\Delta^{-1}\Lambda^{\gamma}(\psi\phi_{\delta})\|_{H^{1}(0,T;H^{-1}(\mathbb{T}^{3}))}\leq
					C\|\Delta^{-1}\Lambda^{\gamma}(\psi\phi_{\delta})\|_{H^{1}(0,T;L^{2}(\mathbb{T}^{3}))}
					\leq C\|\psi\phi_{\delta}\|_{H^{1}(0,T;L^{2}(\mathbb{T}^{3}))}
				\end{align*}
				for $0\leq\gamma\leq1$ and $\psi\in C_{0}^{\infty}(0,T)$.
				%Firstly, we aim to prove
				%\begin{align}
				%  (\mathrm{div}\mathbf{u}^{N},\partial_{t,h}^{+}(\Delta^{-1}\Lambda^{\gamma}(\psi\phi_{\delta})))_{Q_{T}}\rightarrow_{h\rightarrow 0}(\mathrm{div}\mathbf{u}_{\delta},\partial_{t}(\Delta^{-1}\Lambda^{\gamma}(\psi\phi_{\delta})))_{Q_{T}}.\label{dif-limit-1}
				%\end{align}
				%To proceed we write
				%\begin{align}
				%  &(\mathrm{div}\mathbf{u}^{N},\partial_{t,h}^{+}(\psi\Delta^{-1}\Lambda^{\gamma}\phi_{\delta}))_{Q_{T}}-(\mathrm{div}\mathbf{u}_{\delta},\partial_{t}(\psi\Delta^{-1}\Lambda^{\gamma}\phi_{\delta}))_{Q_{T}}
				%  %\nonumber\\&=(\mathrm{div}\mathbf{u}^{N}-\mathrm{div}\mathbf{u}_{\delta},\partial_{t,h}^{+}(\psi\Delta^{-1}\Lambda^{\gamma}\phi_{\delta}))_{Q_{T}}+(\mathrm{div}\mathbf{u}_{\delta},\partial_{t,h}^{+}(\psi\Delta^{-1}\Lambda^{\gamma}\phi_{\delta})-\partial_{t}(\psi\Delta^{-1}\Lambda^{\gamma}\phi_{\delta}))_{Q_{T}}
				%   \nonumber\\&=\big(\mathrm{div}\mathbf{u}^{N},\Delta^{-1}\Lambda^{\gamma}(\partial_{t,h}^{+}(\psi\phi_{\delta}))\big)_{Q_{T}}-\big(\mathrm{div}\mathbf{u}_{\delta},\Delta^{-1}\Lambda^{\gamma}(\partial_{t}(\psi\phi_{\delta}))\big)_{Q_{T}}
				%  \nonumber\\&=\big(\mathrm{div}\mathbf{u}^{N},\Delta^{-1}\Lambda^{\gamma}(\partial_{t,h}^{+}(\psi\phi_{\delta})-\partial_{t}(\psi\phi_{\delta}))\big)_{Q_{T}}+\big(\mathrm{div}\mathbf{u}^N-\mathrm{div}\mathbf{u}_{\delta},\Delta^{-1}\Lambda^{\gamma}(\partial_{t}(\psi\phi_{\delta}))\big)_{Q_{T}}.
				% % \nonumber\\&=(\mathrm{div}\mathbf{u}^{N},\partial_{t,h}^{+}(\psi\Delta^{-1}\Lambda^{\gamma}\phi_{\delta})-\partial_{t}(\psi\Delta^{-1}\Lambda^{\gamma}\phi_{\delta}))_{Q_{T}}
				%%  +(\mathrm{div}\mathbf{u}^{N}-\mathrm{div}\mathbf{u}_{\delta},\partial_{t}(\psi\Delta^{-1}\Lambda^{\gamma}\phi_{\delta}))_{Q_{T}}
				%%   \nonumber\\&=\big(\mathrm{div}\mathbf{u}^{N},\Delta^{-1}\Lambda^{\gamma}(\partial_{t,h}^{+}\phi_{\delta}-\partial_{t}\phi_{\delta})\psi(t+h)\big)
				%%+\big(\mathrm{div}\mathbf{u}^{N},\big(\partial_{t,h}^{+}\psi-\psi'(t)\big)\Delta^{-1}\Lambda^{\gamma}\phi_{\delta}\big)
				%%\nonumber\\&\quad+\big(\mathrm{div}\mathbf{u}^{N},\Delta^{-1}\Lambda^{\gamma}\big(\partial_{t}\phi_{\delta}\big)(\psi(t+h)-\psi(t))\big)+(\mathrm{div}\mathbf{u}^{N}-\mathrm{div}\mathbf{u}_{\delta},\partial_{t}(\psi\Delta^{-1}\Lambda^{\gamma}\phi_{\delta}))_{Q_{T}}
				%   \label{dif-estimation-1}
				%\end{align}
				%Due to $\mathbf{u}^{N}\rightharpoonup_{h\rightarrow0}\mathbf{u}_{\delta}$ in $L^{2}(0,T;H^{1}(\mathbb{T}^{3}))$, we easily find the last term in \eqref{dif-estimation-1} vanish as $h\rightarrow 0$. For the first term in  \eqref{dif-estimation-1},
				%we next verify $\partial_{t,h}^{+}(\psi\phi_{\delta})\rightarrow \partial_{t}(\psi\phi_{\delta})$ strongly in $L^{2}(Q_{T})$ as $h\rightarrow 0$. In fact,
				%First part we verify
				%\begin{align*}
				%&|(\mathrm{div}\mathbf{u}^{N},\partial_{t,h}^{+}(\psi\Delta^{-1}\Lambda^{\gamma}\phi_{\delta})-\partial_{t}(\psi\Delta^{-1}\Lambda^{\gamma}\phi_{\delta}))_{Q_{T}}|\rightarrow_{h\rightarrow 0}0
				%%\\&\leq\|\mathrm{div}\mathbf{u}^{N}\|_{L^{2}(0,T;L^{2}(\mathbb{T}^{3}))}\|\partial_{t,h}^{+}(\psi\Delta^{-1}\Lambda^{\gamma}\phi_{\delta})-\partial_{t}(\psi\Delta^{-1}\Lambda^{\gamma}\phi_{\delta})\|_{L^{2}(0,T;L^{2}(\mathbb{T}^{3}))}
				%%\\&=\|\mathrm{div}\mathbf{u}^{N}\|_{L^{2}(0,T;L^{2}(\mathbb{T}^{3}))}\|\partial_{t,h}^{+}(\psi\Delta^{-1}\Lambda^{\gamma}\phi_{\delta})-\partial_{t}(\psi\Delta^{-1}\Lambda^{\gamma}\phi_{\delta})\|_{L^{2}(0,T;L^{2}(\mathbb{T}^{3}))}
				%\end{align*}
				%
				%\begin{align*}
				%&\partial_{t,h}^{+}(\psi\Delta^{-1}\Lambda^{\gamma}\phi_{\delta})-\partial_{t}(\psi\Delta^{-1}\Lambda^{\gamma}\phi_{\delta})
				%\\&=\Delta^{-1}\Lambda^{\gamma}(\partial_{t,h}^{+}\phi_{\delta})\psi(t+h)+\partial_{t,h}^{+}\psi\Delta^{-1}\Lambda^{\gamma}(\phi_{\delta})-\psi'(t)\Delta^{-1}\Lambda^{\gamma}\phi_{\delta}-\psi\Delta^{-1}\Lambda^{\gamma}(\partial_{t}\phi_{\delta})
				%\\&=\Delta^{-1}\Lambda^{\gamma}(\partial_{t,h}^{+}\phi_{\delta}-\partial_{t}\phi_{\delta})\psi(t+h)
				%+\big(\partial_{t,h}^{+}\psi-\psi'(t)\big)\Delta^{-1}\Lambda^{\gamma}\phi_{\delta}+\Delta^{-1}\Lambda^{\gamma}\big(\partial_{t}\phi_{\delta}\big)(\psi(t+h)-\psi(t))
				%\end{align*}
				%Therefore
				%\begin{align*}
				%&(\mathrm{div}\mathbf{u}^{N},\partial_{t,h}^{+}(\psi\Delta^{-1}\Lambda^{\gamma}\phi_{\delta})-\partial_{t}(\psi\Delta^{-1}\Lambda^{\gamma}\phi_{\delta}))_{Q_{T}}
				%\\&=(\mathrm{div}\mathbf{u}^{N},\Delta^{-1}\Lambda^{\gamma}(\partial_{t,h}^{+}\phi_{\delta}-\partial_{t}\phi_{\delta})\psi(t+h))
				%+(\mathrm{div}\mathbf{u}^{N},\big(\partial_{t,h}^{+}\psi-\psi'(t)\big)\Delta^{-1}\Lambda^{\gamma}\phi_{\delta})
				%\\&+\big(\mathrm{div}\mathbf{u}^{N},\Delta^{-1}\Lambda^{\gamma}\big(\partial_{t}\phi_{\delta}\big)(\psi(t+h)-\psi(t))\big)
				%\end{align*}
				%Estimate above one by one
				%\begin{align*}
				%&|(\mathrm{div}\mathbf{u}^{N},\Delta^{-1}\Lambda^{\gamma}(\partial_{t,h}^{+}\phi_{\delta}-\partial_{t}\phi_{\delta})\psi(t+h))|
				%\\&\leq\|\mathrm{div}\mathbf{u}^{N}\|_{L^{2}(0,T;L^{2}(\mathbb{T}^{3}))}
				%\|\psi(t+h)(\Delta^{-1}\Lambda^{\gamma}(\partial_{t,h}^{+}\phi_{\delta}-\partial_{t}\phi_{\delta}))\|_{L^{2}(0,T;L^{2}(\mathbb{T}^{3}))}
				%\\&\leq C\|\psi(t+h)(\partial_{t,h}^{+}\phi_{\delta}-\partial_{t}\phi_{\delta})\|_{L^{2}(0,T;L^{2}(\mathbb{T}^{3}))}
				%\leq C\|(\partial_{t,h}^{+}\phi_{\delta}-\partial_{t}\phi_{\delta})\|_{L^{2}(0,T-h;L^{2}(\mathbb{T}^{3}))}\rightarrow_{h\rightarrow 0}0,
				%\end{align*}
				%We can first show that
				%$\|\partial_{t,h}^{+}\phi_{\delta}\|_{L^{2}(0,T-h;L^{2}(\mathbb{T}^{3}))}\leq \|\partial_{t}\phi_{\delta}\|_{L^{2}(0,T;L^{2}(\mathbb{T}^{3}))}$.
				%noting that
				%\begin{align*}
				%&\phi_{\delta}(t+h)-\phi_{\delta}(t)=h\int_0^{1} \partial_{t}\phi_{\delta}(t+\tau h)\mathrm{d}\tau,
				%\end{align*}
				%\begin{align*}
				%\partial_{t,h}^{+}(\psi\phi_{\delta})=\int_0^{1} \partial_{t}(\psi\phi_{\delta})(t+\tau h)\mathrm{d}\tau
				%\end{align*}
				%which leads to
				%\begin{align*}
				%&|\partial_{t,h}^{+}(\psi\phi_{\delta})|^{2}=\bigg|\int_0^{1} \partial_{t}(\psi\phi_{\delta})(t+\tau h)\mathrm{d}\tau\bigg|^{2}\leq \int_0^{1} |\partial_{t}(\psi\phi_{\delta})(t+\tau h)|^{2}\mathrm{d}\tau.
				%\end{align*}
				%Consequently, we have
				%\begin{align}
				%&\int_{\mathbb{T}^{3}}\int_0^{T}|\partial_{t,h}^{+}(\psi\phi_{\delta})|^{2}\mathrm{d}t\mathrm{d}x\leq \int_{\mathbb{T}^{3}}\int_0^{T}\int_0^{1} |\partial_{t}(\psi\phi_{\delta})(t+\tau h)|^{2}\mathrm{d}\tau\mathrm{d}t\mathrm{d}x\leq  \int_{\mathbb{T}^{3}}\int_{0}^{T}|\partial_{t}\phi_{\delta}(t )|^{2}\mathrm{d}t\mathrm{d}x.\label{dif-limit-2}
				%\end{align}
				%%integrating over $\mathbb{T}^{3}$, we have
				%%\begin{align*}
				%%&\int_{\mathbb{T}^{3}}\int_0^{T-h}|\partial_{t,h}^{+}\phi_{\delta}|^{2}\mathrm{d}t\mathrm{d}x\leq C \int_{\mathbb{T}^{3}}\int_{0}^{T}|\partial_{t}\phi_{\delta}(t )|^{2}\mathrm{d}t\mathrm{d}x,
				%%\end{align*}
				%Combining this  with the fact that $\partial_{t,h}^{+}(\psi\phi_{\delta})$ converges pointwise to  $ \partial_{t}(\psi\phi_{\delta})$  in $Q_T$  as $h\rightarrow 0$, it follows that $\partial_{t,h}^{+}\phi_{\delta}\rightharpoonup_{h\rightarrow 0} \partial_{t}\phi_{\delta}$ in $L^{2}(Q_{T})$ and
				%\begin{align*}
				%\int_{0}^{T}|\partial_{t}\phi_{\delta}(t )|^{2}\mathrm{d}t\mathrm{d}x&\leq \liminf\limits_{h\rightarrow 0}\int_{\mathbb{T}^{3}}\int_0^{T}|\partial_{t,h}^{+}(\psi\phi_{\delta})|^{2}\mathrm{d}t\mathrm{d}x
				%\nonumber\\ &\leq \limsup\limits_{h\rightarrow 0}\int_{\mathbb{T}^{3}}\int_0^{T}|\partial_{t,h}^{+}(\psi\phi_{\delta})|^{2}\mathrm{d}t\mathrm{d}x\nonumber\\ &\leq \int_{\mathbb{T}^{3}}\int_{0}^{T}|\partial_{t}\phi_{\delta}(t )|^{2}\mathrm{d}t\mathrm{d}x,
				%\end{align*}
				%that is,
				%\begin{align}
				%\lim\limits_{h\rightarrow 0}\int_{\mathbb{T}^{3}}\int_0^{T}|\partial_{t,h}^{+}(\psi\phi_{\delta})|^{2}\mathrm{d}t\mathrm{d}x
				%=\int_{\mathbb{T}^{3}}\int_{0}^{T}|\partial_{t}\phi_{\delta}(t )|^{2}\mathrm{d}t\mathrm{d}x.\label{dif-limit-3}
				%\end{align}
				%We then can get $\partial_{t,h}^{+}(\psi\phi_{\delta})\rightarrow \partial_{t}(\psi\phi_{\delta})$ strongly in $L^{2}(Q_{T})$ as $h\rightarrow 0$.
				%Secondly, we shall prove
				%\begin{align}
				%(\rho^{N}(\rho_{h}^{N})^{-1}\mathbf{u}^{N}\cdot\nabla \mathbf{u}^{N},\psi\nabla\Delta^{-1}\Lambda^{\gamma}\phi_{\delta})_{Q_{T}}\rightarrow_{h\rightarrow 0}(\mathbf{u}_{\delta}\cdot\nabla \mathbf{u}_{\delta},\psi\nabla\Delta^{-1}\Lambda^{\gamma}\phi_{\delta})_{Q_{T}}.\label{dif-limit-4}
				%\end{align}
				%We have
				%\begin{align}
				%&(\rho^{N}(\rho_{h}^{N})^{-1}\mathbf{u}^{N}\cdot\nabla \mathbf{u}^{N},\psi\nabla\Delta^{-1}\Lambda^{\gamma}\phi_{\delta})_{Q_{T}}-(\mathbf{u}_{\delta}\cdot\nabla \mathbf{u}_{\delta},\psi\nabla\Delta^{-1}\Lambda^{\gamma}\phi_{\delta})_{Q_{T}}\nonumber
				%\\&=\big((\rho^{N}(\rho_{h}^{N})^{-1}-1)\mathbf{u}^{N}\cdot\nabla \mathbf{u}^{N},\psi\nabla\Delta^{-1}\Lambda^{\gamma}\phi_{\delta}\big)_{Q_{T}}+
				%\big((\mathbf{u}^{N}-\mathbf{u}_{\delta})\cdot\nabla \mathbf{u}^{N},\psi\nabla\Delta^{-1}\Lambda^{\gamma}\phi_{\delta}\big)_{Q_{T}}\nonumber
				%\\&\quad+\big(\mathbf{u}_{\delta}\cdot\nabla (\mathbf{u}^{N}-\mathbf{u}_{\delta}),\psi\nabla\Delta^{-1}\Lambda^{\gamma}\phi_{\delta}\big)_{Q_{T}}.\label{dif-limit-4-1}
				%\end{align}
				%Thanks to \eqref{N-energy},  we have $\mathbf{u}^{N}\cdot\nabla \mathbf{u}^{N}$ is uniformly bounded in $L^{2}(0,T;L^{1}(\mathbb{T}^{3}))$ which and \eqref{dif-limit-4-1-1} imply
				%\begin{align}
				%&((\rho^{N}(\rho_{h}^{N})^{-1}-1)\mathbf{u}^{N}\cdot\nabla \mathbf{u}^{N},\psi\nabla\Delta^{-1}\Lambda^{\gamma}\phi_{\delta})_{Q_{T}}
				%%\\&\leq \|(\rho^{N}(\rho_{h}^{N})^{-1}-1)\|_{L^{\infty}(0,T;L^{\infty}(\mathbb{T}^{3}))}
				%%\|\mathbf{u}^{N}\|_{L^{\infty}(0,T;L^{2}(\mathbb{T}^{3}))}
				%%\\&\quad\|\nabla\mathbf{u}^{N}\|_{L^{2}(0,T;L^{2}(\mathbb{T}^{3}))}\|\nabla\Delta^{-1}\Lambda^{\gamma}\phi_{\delta}\|_{L^{2}(0,T;L^{\infty}(\mathbb{T}^{3}))}
				%%\\&\leq\|\psi((\rho^{N}(\rho_{h}^{N})^{-1}-1))\|_{L^{\infty}(0,T;L^{\infty}(\mathbb{T}^{3}))}
				%\rightarrow_{h\rightarrow 0}0,\label{dif-limit-4-1-2}
				%\end{align}
				%where we have used  $\rho^{N}\rightarrow \rho_\delta$  in $L^{2}(0,T;H^{s}(\mathbb{T}^{3}))\hookrightarrow L^{2}(0,T;L^{\infty}(\mathbb{T}^{3}))$ as $h\rightarrow 0$.
				%
				%
				%Moreover, since that $\mathbf{u}^{N}\rightarrow \mathbf{u}_{\delta}$ in $L^{2}(0,T;L^{2}(\mathbb{T}^{3}))$(cf.\eqref{Strong convergence 1}) and $\nabla\mathbf{u}^{N}\rightharpoonup \nabla\mathbf{u}_{\delta}$ in $L^{2}(0,T;L^{2}(\mathbb{T}^{3}))$(cf. \eqref{weak convergence 1}) as $h\rightarrow 0$, there hold
				%\begin{align}
				%&((\mathbf{u}^{N}-\mathbf{u}_{\delta})\cdot\nabla \mathbf{u}^{N},\psi\nabla\Delta^{-1}\Lambda^{\gamma}\phi_{\delta})_{Q_{T}}\rightarrow_{h\rightarrow 0}0\label{dif-limit-4-1-3}
				%\end{align}
				%and
				%\begin{align}
				%&(\mathbf{u}_{\delta}\cdot\nabla (\mathbf{u}^{N}-\mathbf{u}_{\delta}),\psi\nabla\Delta^{-1}\Lambda^{\gamma}\phi_{\delta})_{Q_{T}}\rightarrow_{h\rightarrow 0}0.\label{dif-limit-4-1-4}
				%\end{align}
				%By substituting \eqref{dif-limit-4-1-2}-\eqref{dif-limit-4-1-4} into \eqref{dif-limit-4-1} we immediately derive \eqref{dif-limit-4}.
				%
				%Thirdly, we show
				%\begin{align}
				%(\nabla(\rho_{h}^{N})^{-1}S(\phi_{h}^{N},\mathbb{ D}\mathbf{u}^{N}),\psi\nabla\Delta^{-1}\Lambda^{\gamma}\phi_{\delta})_{Q_{T}}\rightarrow_{h\rightarrow 0}(\nabla(\rho_{\delta})^{-1}S(\phi_{\delta},\mathbb{ D}\mathbf{u}_{\delta}),\psi\nabla\Delta^{-1}\Lambda^{\gamma}\phi_{\delta})_{Q_{T}}.\label{dif-limit-5}
				%\end{align}
				%\begin{align}\label{weak1110}
				%  &\nonumber(\mathrm{div}\mathbf{u}^{N},\partial_{t,h}^{+}(\psi\Delta^{-1}\Lambda^{\gamma}\phi_{\delta}))_{Q_{T}}+(\rho^{N}(\rho_{h}^{N})^{-1}\mathbf{u}^{N}\cdot\nabla \mathbf{u}^{N},\psi\nabla\Delta^{-1}\Lambda^{\gamma}\phi_{\delta})_{Q_{T}}
				%  \\&\nonumber\quad+((\rho_{h}^{N})^{-1}S(\phi_{h}^{N},\mathbb{ D}\mathbf{u}^{N}),\nabla(\psi\nabla\Delta^{-1}\Lambda^{\gamma}\phi_{\delta}))_{Q_{T}}
				%  +(\nabla(\rho_{h}^{N})^{-1}S(\phi_{h}^{N},\mathbb{ D}\mathbf{u}^{N}),\psi\nabla\Delta^{-1}\Lambda^{\gamma}\phi_{\delta})_{Q_{T}}
				% \\&\nonumber\quad -\zeta((\rho_{h}^{N})^{-1}\rho^{N}p_{0}^{N},\mathrm{div}(\psi\nabla\Delta^{-1}\Lambda^{\gamma}\phi_{\delta}))_{Q_{T}}
				%  +((\rho_{h}^{N})^{-1}\phi^{N}\nabla\mu_{p}^{0,N},\psi\nabla\Delta^{-1}\Lambda^{\gamma}\phi_{\delta})_{Q_{T}}
				%  \\&\quad-\zeta(\nabla((\rho_{h}^{N})^{-1}\rho^{N}) p_{0}^{N},\psi\nabla\Delta^{-1}\Lambda^{\gamma}\phi_{\delta})_{Q_{T}}
				% +(\frac{\varepsilon\delta}{4\alpha}(\rho_{h}^{N})^{-1}\mathbf{u}^{N}p_{0}^{N},\psi\nabla\Delta^{-1}\Lambda^{\gamma}\phi_{\delta})_{Q_{T}}=0,
				%\end{align}
				%
				%Finally, by using similar arguments, one has
				
				Consequently we arrive at
				\begin{align}\label{weak1111}
					&\nonumber(\mathrm{div}\mathbf{u}_{\delta},\partial_{t}(\Delta^{-1}\Lambda^{\gamma}(\psi\phi_{\delta})))_{Q_{T}}-(\mathbf{u}_{\delta}\cdot\nabla \mathbf{u}_{\delta},\nabla\Delta^{-1}\Lambda^{\gamma}(\psi\phi_{\delta}))_{Q_{T}}
					\\&\nonumber\quad+(\nabla(\rho_{\delta})^{-1}S(\phi_{\delta},\mathbb{ D}\mathbf{u}_{\delta}),\nabla\Delta^{-1}\Lambda^{\gamma}(\psi\phi_{\delta}))_{Q_{T}}+((\rho_{\delta}^{-1}S(\phi_{\delta},\mathbb{ D}\mathbf{u}_{\delta}),\nabla(\nabla\Delta^{-1}\Lambda^{\gamma}(\psi\phi_{\delta})))_{Q_{T}}
					\\&\nonumber\quad -\zeta(p_{0,\delta},\mathrm{div}(\nabla\Delta^{-1}\Lambda^{\gamma}(\psi\phi_{\delta})))_{Q_{T}}
					+((\rho_{\delta})^{-1}\phi_{\delta}\nabla\mu_{p,\delta}^{0},\nabla\Delta^{-1}\Lambda^{\gamma}(\psi\phi_{\delta}))_{Q_{T}}
					\\&\quad
					+(\frac{\varepsilon\delta}{4\alpha}(\rho_{\delta})^{-1}\mathbf{u}_{\delta}p_{0,\delta},\nabla\Delta^{-1}\Lambda^{\gamma}(\psi\phi_{\delta}))_{Q_{T}}=0.
				\end{align}
				
				We  submit  $p_{0,\delta}=\frac{1}{\alpha}(\mu_{p,\delta}^{0}-\Lambda^{2s}\phi_{\delta}-F'(\phi_{\delta})+\bar{\mu}_{\delta})$ into \eqref{weak1111} and then get
				\begin{align}
					& (\zeta\alpha^{-1}\Lambda^{s+\frac{\gamma}{2}}\phi_{\delta},\Lambda^{s+\frac{\gamma}{2}}(\psi\phi_{\delta}))_{Q_{T}}
					\nonumber \\& =%-( \mathrm{div}\mathbf{u}_{\delta},\partial_t\psi\Delta^{-1}\Lambda^{\gamma} \phi_{\delta})_{Q^{T}}
					(\mathbf{u}_{\delta},\nabla\Delta^{-1}\Lambda^{\gamma}\partial_t (\psi\phi_{\delta}))_{Q^{T}}
					+(\mathbf{u}_{\delta}\cdot\nabla
					\mathbf{u}_{\delta},\nabla\Delta^{-1}\Lambda^{\gamma}(\psi\phi_{\delta}))_{Q_{T}}
					\nonumber \\&\quad+((\rho_{\delta}^{-1}S(\phi_{\delta},\mathbb{ D}\mathbf{u}_{\delta}),\nabla(\nabla\Delta^{-1}\Lambda^{\gamma}(\psi\phi_{\delta})))_{Q_{T}}
					\nonumber \\&\quad+(\nabla(\rho_{\delta})^{-1}S(\phi_{\delta},\mathbb{ D}\mathbf{u}_{\delta}),\nabla\Delta^{-1}\Lambda^{\gamma}(\psi\phi_{\delta}))_{Q_{T}}
					\nonumber\\&\quad +(\phi_{\delta}\rho_{\delta}^{-1} \nabla \mu_{p,\delta},\nabla\Delta^{-1}\Lambda^{\gamma}(\psi\phi_{\delta}))_{Q_{T}}
					\nonumber\\& \quad+\bigg(\frac{\varepsilon\delta}{4\alpha\rho_{\delta}}  p_{0,\delta} \mathbf{u}_{\delta},\nabla\Delta^{-1}\Lambda^{\gamma}(\psi\phi_{\delta})\bigg)_{Q_{T}}
					\nonumber \\&\quad-\frac{\zeta}{\alpha}(\nabla(\mu_{p,\delta}^{0}-F'(\phi_{\delta})),\nabla\Delta^{-1}\Lambda^{\gamma}(\psi\phi_{\delta}))_{Q_{T}}
					\nonumber \\&\eqqcolon\sum_{i=1}^{7}J_{i}.\label{lemh55-5}
				\end{align}
				
				
				
				
				%By basic energy estimate \eqref{ConEnL2} and H\"{o}lder inequality, we have for $\gamma\leq s+2$
				% \begin{align}
					%    \label{lemh55-6}
					%   \nonumber  |J_{1}|
					%    &\lesssim  \|\mathrm{div}\mathbf{u}_{\delta}\|_{L^{2}(0,T;L^{2}(\mathbb{T}^{3}))}\|\partial_t\psi\|_{L^{2}(0,T)}\|\Delta^{-1}\Lambda^{\gamma} \phi_{\delta}\|_{L^{\infty}(0,T;L^{2}(\mathbb{T}^{3}))}
					%   \\& \lesssim \|\mathbf{u}_{\delta}\|_{L^{2}(0,T;H^{1}(\mathbb{T}^{3}))}\|\partial_t\psi\|_{L^{2}(0,T)}\| \phi_{\delta}\|_{L^{\infty}(0,T;H^{\gamma-2}(\mathbb{T}^{3}))}
					%   \lesssim\|\partial_t\psi\|_{L^{2}(0,T)}.
					%  \end{align}
				Furthermore, in view of \eqref{model4-3} and \eqref{ConEnL2} we can find
				\begin{align}\|\partial_{t}\phi_{\delta}\|_{L^{2}(0,T;L^{2}(\mathbb{T}^{3}))}&=\|\mathrm{div}(\phi_{\delta}  \mathbf{u}_{\delta})-\Delta\mu_{p,\delta}^{0}\|_{L^{2}(0,T;L^{2}(\mathbb{T}^{3}))}
					\nonumber\\&=\|\phi_{\delta}\mathrm{div}\mathbf{u}_{\delta}+\nabla\phi_{\delta}\cdot  \mathbf{u}_{\delta}-\alpha^{-1}(\mathrm{div}  \mathbf{u}_{\delta} +\delta p_{0,\delta})\|_{L^{2}(0,T;L^{2}(\mathbb{T}^{3}))}
					\nonumber\\&\lesssim \|\phi_{\delta}\|_{L^{\infty}(Q_T))}
					\|\mathrm{div}\mathbf{u}_{\delta}\|_{L^{2}(Q_T)}+\|\nabla\phi_{\delta}\|_{L^{\infty}(0,T;L^{3}(\mathbb{T}^{3}))}
					\|\mathbf{u}_{\delta}\|_{L^{2}(0,T;L^{6}(\mathbb{T}^{3}))}
					\nonumber\\& \quad+\alpha^{-1}\|\mathrm{div}  \mathbf{u}_{\delta}\|_{L^{2}(0,T;L^{2}(\mathbb{T}^{3}))}
					+\alpha^{-1}\|\delta p_{0,\delta}\|_{L^{2}(0,T;L^{2}(\mathbb{T}^{3}))}
					\nonumber\\&\lesssim \alpha^{-1}                  \label{uniform phi-alpha}
				\end{align}  which and \eqref{ConEnL2} imply
				\begin{align}
					\label{lemh55-7}
					\nonumber  |J_{1}|&
					\lesssim \|\mathbf{u}_{\delta}\|_{L^{\infty}(0,T;L^{2}(\mathbb{T}^{3}))}\|\nabla\Delta^{-1}\Lambda^{\gamma}\partial_t (\psi\phi_{\delta})\|_{L^{1}(0,T;L^{2}(\mathbb{T}^{3}))}
					% \\& \nonumber\lesssim \|\mathbf{u}_{\delta}\|_{L^{2}(0,T;H^{1}(\mathbb{T}^{3}))}\|\psi\|_{L^{\infty}(0,T)}\|\partial_t (\psi\phi_{\delta})\|_{L^{2}(0,T;H^{\gamma-2}(\mathbb{T}^{3}))}
					\\& \nonumber\lesssim (\|\psi\|_{L^{\infty}(0,T)}\|\partial_t \phi_{\delta}\|_{L^{2}(0,T;L^{2}(\mathbb{T}^{3}))}+\|\partial_t\psi\|_{L^{1}(0,T)}\|\phi_{\delta}\|_{L^{\infty}(0,T;L^{2}(\mathbb{T}^{3}))})
					\\&\lesssim (\frac{1}{\alpha}\|\psi\|_{L^{\infty}(0,T)}+\|\partial_t\psi\|_{L^{1}(0,T)})
				\end{align}
				for $\gamma\leq 1$.
				Similarly, using the fact that $H^{s-\gamma+1}\hookrightarrow L^{\infty}$ for $\gamma\leq 1<s-\frac{1}{2}$, we derive
				\begin{align}
					\label{lemh55-8}
					\nonumber  |J_{2}|
					& \lesssim \|\mathbf{u}_{\delta}\|_{L^{\infty}(0,T;L^{2}(\mathbb{T}^{3}))}\|\nabla\mathbf{u}_{\delta}\|_{L^{2}(0,T;L^{2}(\mathbb{T}^{3}))}\|\nabla\Delta^{-1}\Lambda^{\gamma} (\psi\phi_{\delta})\|_{L^{2}(0,T;L^{\infty}(\mathbb{T}^{3}))}
					\\&\nonumber\lesssim \|\mathbf{u}_{\delta}\|_{L^{\infty}(0,T;L^{2}(\mathbb{T}^{3}))}\|\psi\|_{L^{2}(0,T)}\|\mathbf{u}_{\delta}\|_{L^{2}(0,T;H^{1}(\mathbb{T}^{3}))}\|\nabla\Delta^{-1}\Lambda^{\gamma} (\psi\phi_{\delta})\|_{L^{2}(0,T;H^{s-\gamma+1}(\mathbb{T}^{3}))}
					\\&\nonumber\lesssim \|\mathbf{u}_{\delta}\|_{L^{\infty}(0,T;L^{2}(\mathbb{T}^{3}))}\|\psi\|_{L^{\infty}(0,T)}\|\mathbf{u}_{\delta}\|_{L^{2}(0,T;H^{1}(\mathbb{T}^{3}))}\| \phi_{\delta}\|_{L^{\infty}(0,T;H^{s}(\mathbb{T}^{3}))}
					\\&\lesssim \|\psi\|_{L^{\infty}(0,T)}.
				\end{align}
				
				
				
				For $J_{3}$ and $J_{4}$,  it follows from $H^{s}(\mathbb{T}^{3})\hookrightarrow W^{1,3}(\mathbb{T}^{3})(s>\frac{3}{2})$, $H^{1}(\mathbb{T}^{3})\hookrightarrow L^{6}(\mathbb{T}^{3})$ and H\"{o}lder inequality  that
				\begin{align}
					\label{lemh55-9}
					\nonumber |J_{3}|+|J_{4}|& \lesssim |(\nabla\rho_{\delta}^{-1}S(\phi_{\delta},\mathbb{ D}\mathbf{u}_{\delta}),\nabla\Delta^{-1}\Lambda^{\gamma}(\psi\phi_{\delta}))_{Q_{T}}|+|(\rho_{\delta}^{-1}S(\phi_{\delta},\mathbb{ D}\mathbf{u}_{\delta}),\nabla^{2}\Delta^{-1}\Lambda^{\gamma}(\psi\phi_{\delta}))_{Q_{T}}|
					\\&\nonumber\lesssim \|\psi\|_{L^{2}(0,T)}\|\nabla\phi_{\delta}\|_{L^{\infty}(0,T;L^{3}(\mathbb{T}^{3}))}\|\mathbb{ D}\mathbf{u}_{\delta}\|_{L^{2}(0,T;L^{2}(\mathbb{T}^{3}))}\|\nabla\Delta^{-1}\Lambda^{\gamma} \phi_{\delta}\|_{L^{\infty}(0,T;L^{6}(\mathbb{T}^{3}))}
					\\&\nonumber \quad+\|\psi\|_{L^{2}(0,T)}\|\mathbb{ D}\mathbf{u}_{\delta}\|_{L^{2}(0,T;L^{2}(\mathbb{T}^{3}))}\|\nabla^{2}\Delta^{-1}\Lambda^{\gamma} \phi_{\delta}\|_{L^{\infty}(0,T;L^{2}(\mathbb{T}^{3}))}
					\\&\nonumber\lesssim \|\psi\|_{L^{\infty}(0,T)}\|\phi_{\delta}\|_{L^{\infty}(0,T;H^{s}(\mathbb{T}^{3}))}\|\mathbf{u}_{\delta}\|_{L^{2}(0,T;H^{1}(\mathbb{T}^{3}))}\| \phi_{\delta}\|_{L^{\infty}(0,T;H^{\gamma}(\mathbb{T}^{3}))}
					\\&\nonumber \quad+\|\psi\|_{L^{\infty}(0,T)}\|\mathbf{u}_{\delta}\|_{L^{2}(0,T;H^{1}(\mathbb{T}^{3}))}\|\phi_{\delta}\|_{L^{\infty}(0,T;H^{\gamma}(\mathbb{T}^{3}))}
					\\&\lesssim\|\psi\|_{L^{\infty}(0,T)}.
				\end{align}
				
				By using the similar argument we can easily deduce
				\begin{align}
					\label{lemh55-10}
					|J_{5}|\lesssim \|\psi\|_{L^{\infty}(0,T)},\ |J_{6}|\lesssim \|\psi\|_{L^{\infty}(0,T)},\ |J_{7}|\lesssim \frac{1}{\alpha}\|\psi\|_{L^{\infty}(0,T)}.
				\end{align}
				%By H\"{o}lder inequality and the fact that $\alpha=-\frac{\varepsilon}{2+\varepsilon}$ and $0 < \abs{\varepsilon} < 1$ is small,
				% \begin{align}
					%   &  \label{lemh55-11}
					%  |J_{6}|
					% \lesssim \|\psi\|_{L^{2}(0,T)}\|\delta p_{0,\delta}\|_{L^{2}(0,T;L^{2}(\mathbb{T}^{3}))}\| \phi\|_{L^{\infty}(0,T;H^{\gamma-1}(\mathbb{T}^{3}))}\lesssim \|\psi\|_{L^{\infty}(0,T)}.
					%  \end{align}
				% The last integral can be estimated by the H\"{o}lder's inequality, Assumption \ref{ass:main} and $\gamma\leq s+1$, as follows
				% \begin{align}
					%     \label{lemh55-12}
					% \nonumber |J_{7}|
					% & \lesssim \frac{\zeta}{\alpha}\|\nabla(\mu_{p,\delta}^{0}-F'(\phi_{\delta}))\|_{L^{2}(0,T;L^{2}(\mathbb{T}^{3}))}\|\nabla\Delta^{-1}\Lambda^{\gamma} (\psi\phi_{\delta})\|_{L^{\infty}(0,T;L^{2}(\mathbb{T}^{3}))}
					%\\& \nonumber\lesssim \frac{1}{\alpha}(\|\nabla\mu_{p,\delta}^{0}\|_{L^{2}(0,T;L^{2}(\mathbb{T}^{3}))}+\|F''\nabla\phi\|_{L^{2}(0,T;L^{2}(\mathbb{T}^{3}))})\|\nabla\Delta^{-1}\Lambda^{\gamma} \phi_{\delta}\|_{L^{\infty}(0,T;L^{2}(\mathbb{T}^{3}))}\|\psi\|_{L^{2}(0,T)}
					%\\& \nonumber\lesssim \frac{1}{\alpha}(\|\nabla\mu_{p,\delta}^{0}\|_{L^{2}(0,T;L^{2}(\mathbb{T}^{3}))}+\|F''\|_{L^{\infty}(0,T;L^{\infty}(\mathbb{T}^{3}))}\|\phi_{\delta}\|_{L^{2}(0,T;H^{1}(\mathbb{T}^{3}))})
					%\\& \nonumber\quad\times\| \phi_{\delta}\|_{L^{\infty}(0,T;H^{\gamma-1}(\mathbb{T}^{3}))}\|\psi\|_{L^{2}(0,T)}
					% \\& \lesssim \frac{1}{\alpha}\|\psi\|_{L^{\infty}(0,T)}.
					%  \end{align}
				
				
				Combining \eqref{lemh55-7}--\eqref{lemh55-10} with \eqref{lemh55-5}, we show
				\begin{align}
					\label{lemh55-121}
					&
					(\Lambda^{s+\frac{\gamma}{2}}\phi_\delta,\psi\Lambda^{s+\frac{\gamma}{2}}\phi_\delta)_{Q_{T}}
					\lesssim \alpha\|\partial_t\psi\|_{L^{1}(0,T)}+\|\psi\|_{L^{\infty}(0,T)}.
				\end{align}
				Choosing  $\psi=\psi_{m}$ where
				\begin{equation}
					\label{cut-off function}
					\begin{aligned}
						& \psi_{m}\in C_{0}^{\infty}(0,T),  0\leq\psi_{m}\leq1,
						\\& \psi_{m}(t)=1, t\in\left[\frac{1}{m},T-\frac{1}{m}\right], |\psi'|\leq 2m,\ m\in \mathbb{N},
					\end{aligned}
				\end{equation}
				and then sending  $ m \to \infty $,  one derives $\|\phi_{\delta}\|_{L^{2}(0,T;H^{s+\frac{\gamma}{2}}(\mathbb{T}^{3}))}\leq C$. The proof is completed.
			\end{proof}

\end{document}
